\documentclass{report}
\usepackage{amsmath,amssymb}
\setlength{\parindent}{0mm}
\setlength{\parskip}{1em}
\begin{document}
\begin{center}
\rule{6in}{1pt} \
{\large Stephanie Friedhoff \\
{\bf Multigrid approaches for parallel-in-time integration}}

KU Leuven - University of Leuven \\ Department of Computer Science \\ Celestijnenlaan 200A \\ 3001 Leuven-Heverlee (Belgium)
\\
{\tt stephanie.friedhoff@alumni.tufts.edu}\\
Stefan Vandewalle\end{center}

Multigrid and related multilevel methods are the approaches of choice for
solving the linear systems that result from the discretisation of a wide
class of PDEs. One feasible approach for doing effective parallel-in-time
integration is also through some kind of multilevel method. Multigrid
algorithms for parabolic problems that allow temporal parallelism either
employ a semicoarsening strategy or an adaptive parameter-dependent
coarsening strategy. As shown in [Vandewalle and Van de Velde, Space-time
concurrent multigrid waveform relaxation, Ann. Numer. Math., 1 (1994),
pp. 347-360], the multigrid waveform relaxation method which uses
coarsening only in the spatial dimension also permits time parallelism.
The other semicoarsening approach, i.e., coarsening only in the time
dimension is the basis of the multigrid-reduction-in-time (MGRIT)
algorithm [Falgout et. al., SIAM J. Sci. Comput., 36 (2014), pp.
C635-C661]. In space-time multigrid as presented in [Horton and
Vandewalle, A space-time multigrid method for parabolic PDEs, SIAM J.
Sci. Comput., 16 (1995), pp. 848-864], where time is simply another
dimension in the grid, either semicoarsening in space or in time
depending on the anisotropy of the discretisation stencil at each grid
level is employed.

In this talk, we compare different approaches for integrating multigrid
concepts into time-evolution algorithms.


\end{document}
